%=============================================================================
%  Chapitre 1 — Introduction Générale
%=============================================================================
\chapter{Introduction Générale}

\section{Contexte et Motivation}

Les maisons connectées ont évolué d'écosystèmes de gadgets isolés vers des systèmes
cyber-physiques intégrés où la captation, l'actionnement, l'analyse et l'interaction
utilisateur sont étroitement couplés. Selon Statista (2024), le nombre de dispositifs
de maison intelligente en service dans le monde a dépassé 1,3 milliard d'unités, et le
marché mondial de la maison intelligente est projeté à 622 milliards de dollars d'ici 2026.

Malgré cette croissance, la plupart des solutions disponibles restent fragmentées : une
application mobile pour les caméras, une autre pour l'éclairage, et des outils distincts
pour les dispositifs de sécurité. Cette fragmentation crée des inefficacités opérationnelles,
affaiblit la posture de sécurité et génère des frictions pour les utilisateurs quotidiens.

La \textbf{Plateforme IoT Maison Intelligente} répond à cette lacune avec une solution
unifiée : une plateforme unique de qualité production qui intègre la gestion des appareils,
les canaux de commande et de télémétrie en temps réel, la gouvernance du foyer, les flux
d'urgence sécuritaire, un assistant IA, et des scénarios d'automatisation extensibles ---
le tout accessible depuis une interface web responsive.

\section{Cadre du Projet}

Ce projet a été développé en tant que Projet de Fin d'Études (PFE) à
\textbf{ISET Siliana --- Institut Supérieur des Etudes Technologiques de Siliana}, année
académique 2025--2026, en vue de l'obtention de la Licence Appliquée en Informatique.

Le projet est organisé comme un \textbf{monorepo} contenant trois artefacts déployables :

\begin{itemize}[leftmargin=1.5cm,itemsep=6pt]
  \item \textbf{Service backend} --- API REST Node.js/Express, courtier MQTT Aedes intégré,
        passerelle Socket.IO, et couche de persistance MongoDB.
  \item \textbf{Client frontend} --- Application monopage React~19/Vite avec routage
        conscient des rôles, synchronisation d'état en temps réel, plan d'étage 3D
        isométrique, et interface de chat IA.
  \item \textbf{Firmware ESP32} --- Sketch Arduino sur microcontrôleur ESP32-D intégrant
        la détection de température/humidité (DHT11), la détection de gaz (MQ6), le
        contrôle de LEDs RGB, et l'actionnement de servomoteurs.
\end{itemize}

\section{Problématique}

Le problème d'ingénierie central traité par ce projet est la conception et la mise en
œuvre d'une \textbf{plateforme multi-locataires de maison intelligente scalable} qui
satisfait simultanément les contraintes suivantes :

\begin{enumerate}[leftmargin=1.5cm,itemsep=6pt]
  \item \textbf{Réactivité temps réel} pour les événements utilisateur et la télémétrie
        IoT avec une latence bout en bout inférieure à 500~ms.
  \item \textbf{Séparation claire des rôles} entre les comptes Agence, Administrateur
        de maison, Résident et SuperAdmin avec une gestion granulaire des permissions.
  \item \textbf{Comportement critique de sécurité} pour les incidents gazeux avec des
        transitions d'état d'urgence déterministes, une actuation matérielle et une
        notification multi-canal.
  \item \textbf{Architecture maintenable} supportant l'extension progressive des
        fonctionnalités sans rompre les contrats existants.
  \item \textbf{Intégration IA accessible} permettant le contrôle des appareils en
        langage naturel via une API de modèle de langage à faible latence.
\end{enumerate}

\section{Objectifs du Projet}

Les objectifs du projet, mappés à des livrables mesurables, sont les suivants :

\begin{center}
\begin{longtable}{p{0.08\textwidth} p{0.48\textwidth} p{0.34\textwidth}}
\caption{Objectifs et Livrables du Projet}\label{tab:objectifs}\\
\toprule
\textbf{ID} & \textbf{Objectif} & \textbf{Livrable} \\
\midrule
\endfirsthead
\toprule
\textbf{ID} & \textbf{Objectif} & \textbf{Livrable} \\
\midrule
\endhead
\bottomrule
\endfoot
O1  & Construire un système d'authentification multi-rôles complet et sécurisé & API Auth, middleware JWT, OAuth, 2FA \\
O2  & Activer les flux CRUD et de commande pour les appareils connectés & API Appareils, pipeline de publication MQTT \\
O3  & Intégrer un courtier MQTT pour la messagerie IoT & Courtier Aedes, routage de topics, client ESP32 \\
O4  & Livrer la synchronisation d'état en temps réel aux clients web & Passerelle Socket.IO, catalogue d'événements \\
O5  & Implémenter la surveillance d'urgence gaz avec réponse matérielle & Service Gas Monitor, workflow d'urgence \\
O6  & Fournir un flux de demande de permission pour les résidents & API Permissions, système de notification \\
O7  & Construire un plan d'étage 3D avec contrôle interactif des appareils & Composant FloorPlan, overlay SVG \\
O8  & Développer un assistant IA (Melo) avec contrôle en langage naturel & API Companion, interface chat Melo \\
O9  & Implémenter l'automatisation par scénarios temporels & API Scénarios, planificateur node-cron \\
O10 & Créer un tableau de bord de surveillance de la consommation énergétique & API Énergie, visualisation recharts \\
O11 & Déployer le système complet avec runbook opérationnel documenté & Config déploiement, annexe runbook \\
\end{longtable}
\end{center}

\section{Approche Méthodologique --- Scrum}

La mise en œuvre a suivi le cadre agile \textbf{Scrum}, choisi pour ses cycles de livraison
itératifs, l'intégration continue du feedback et la gestion transparente du backlog.
Les rôles Scrum ont été mappés au contexte académique :

\begin{itemize}[leftmargin=1.5cm,itemsep=6pt]
  \item \textbf{Product Owner} --- Encadrant académique (validation des exigences fonctionnelles).
  \item \textbf{Scrum Master} --- Auteur (facilitation des cérémonies, suivi de la vélocité).
  \item \textbf{Équipe de développement} --- Auteur (développement full-stack et firmware).
\end{itemize}

\noindent Le projet a été livré en \textbf{six sprints} de deux semaines chacun, comme
résumé dans le Tableau~\ref{tab:sprint_resume}.

\begin{center}
\begin{longtable}{c p{0.26\textwidth} p{0.38\textwidth} c}
\caption{Résumé des Sprints}\label{tab:sprint_resume}\\
\toprule
\textbf{Sprint} & \textbf{Thème} & \textbf{Livrables Principaux} & \textbf{SP} \\
\midrule
\endfirsthead
\toprule
\textbf{Sprint} & \textbf{Thème} & \textbf{Livrables Principaux} & \textbf{SP} \\
\midrule
\endhead
\bottomrule
\endfoot
1 & Authentification et Gestion des Utilisateurs & Flux d'inscription, JWT, OAuth, 2FA, vérification email & 34 \\
2 & Gestion des Appareils et MQTT              & CRUD Appareils, courtier MQTT, firmware ESP32           & 39 \\
3 & Plan d'Étage et Synchronisation TR         & Plan 3D, scènes/ambiances, passerelle Socket.IO         & 37 \\
4 & Systèmes de Sécurité                       & Moniteur gaz, DHT11, workflow d'urgence, vanne          & 36 \\
5 & Assistant IA et Automatisation             & Melo IA, scénarios, surveillance énergétique            & 34 \\
6 & Portail Agence et Audit                    & Hub conciergerie, journaux d'audit, permissions         & 37 \\
\midrule
  & \textbf{Total}                             &                                                          & \textbf{217} \\
\bottomrule
\end{longtable}
\end{center}

\section{Contributions Scientifiques et Techniques}

Ce projet contribue au domaine de l'IoT et du génie logiciel web en plusieurs dimensions :

\begin{itemize}[leftmargin=1.5cm,itemsep=6pt]
  \item Une \textbf{architecture pratique} combinant APIs REST, courtage MQTT et WebSockets
        dans un service Node.js cohérent.
  \item Un \textbf{modèle de gouvernance des rôles} adaptable aux contextes résidentiels
        et de gestion de propriétés multi-unités.
  \item Une \textbf{boucle de contrôle sécuritaire} liant les seuils de télémétrie PPM
        aux commandes matérielles et aux notifications multi-canal.
  \item Une \textbf{interface en langage naturel} alimentée par LLM, contrainte par le
        contrôle d'accès basé sur les rôles.
  \item Un \textbf{plan d'artefacts Scrum} applicable aux projets IoT similaires en
        contexte académique.
\end{itemize}

\section{Organisation du Rapport}

Ce rapport est structuré pour guider le lecteur des aspects conceptuels aux détails
d'implémentation et opérationnels :

\begin{description}[leftmargin=1.5cm,itemsep=6pt]
  \item[Chapitres 1--3 :] Contexte du projet, analyse des besoins, état de l'art et
        évaluation technologique.
  \item[Chapitre 4 :] Architecture système et modélisation UML (déploiement, composants,
        cas d'utilisation, classes et diagrammes de séquence).
  \item[Chapitres 5--8 :] Implémentation sprint par sprint couvrant la base de données,
        les APIs backend, l'interface frontend et l'intégration IoT.
  \item[Chapitres 9--11 :] Évaluation sécuritaire, stratégie de test, pipeline de
        déploiement et gestion de projet.
  \item[Chapitre 12 :] Conclusion, rétrospective et feuille de route future.
  \item[Annexes (13--16) :] Référence API complète, inventaire frontend, runbook
        d'installation et matrice de traçabilité.
\end{description}
